\documentclass[12pt]{article}
\usepackage[utf8]{inputenc}
\usepackage[T1]{fontenc}
\usepackage{hyperref}
\title{Purum Language Resources for K-12 Curriculum Development}
\author{Compiled by Claude Sonnet 5.5 (claude-sonnet-5-5) for Andr\'as Kornai}
\date{}
\begin{document}
\maketitle

\section{Introduction}
Purum (ISO 639-3: pub; Glottocode puru1266) is a Tibeto-Burman language of Senapati district, Manipur, India, classified by some as Northwestern (``Old Kuki'') Kuki-Chin. Speakers regard themselves as ethnic Naga. Sources cite about 500 speakers (2001 census) and UNESCO lists it as critically endangered; Meitei is used as a second language. No teaching materials were found.

\section{Primary Educational Resources}
None found.

\section{Linguistic and Reference Materials}
\begin{itemize}
\item Purum language (overview and classification): \url{https://en.wikipedia.org/wiki/Purum_language}
\item Purum people (community background): \url{https://en.wikipedia.org/wiki/Purum_people}
\item Purum word list in the ASJP database: \url{https://asjp.clld.org/languages/PURUM}
\item Purum ethnographic overview (Everyculture): \url{https://www.everyculture.com/South-Asia/Purum-Orientation.html}
\item North East Indian Linguistics (front matter of Hyslop et al., 2011): \url{https://researchonline.jcu.edu.au/17240/11/17240_Hyslop_et_al_2011_Front_Pages.pdf}
\item Wikipedia cites N. Gopendro Singh's 2003 Manipur University doctoral dissertation \emph{A descriptive grammar of Purum}; I did not find a URL for it.
\end{itemize}

\section{Religious and Community Resources}
None found.

\section{Digital and Audio Resources}
None found.

\section{Curriculum Development Context}
Purum is very small and the available material is descriptive only. Curriculum work would depend on community documentation and on the unpublished or hard-to-reach grammar mentioned above. Meitei is the common second language in the area.
\end{document}